%-----------------------------------------------------------------------------------
%
%  Plotting with matplotlib (the PyBridge)
%
%-----------------------------------------------------------------------------------

\chapter{Plotting with matplotlib}
\label{pybridge}

FEMM's scripting engine can drive {\tt matplotlib}, so a script can produce
a publication-quality figure directly from a solution:

\begin{verbatim}
    import matplotlib.pyplot as plt

    fig, ax = plt.subplots()
    ax.plot([0, 1, 2], [0, 1, 4], 'r-')
    ax.set_xlabel("x")
    plt.show()
\end{verbatim}

A real, interactive window appears --- with the usual zoom, pan and save
toolbar --- and it stays alive after the script ends, so you can keep working
with it.

FEMM does not contain {\tt matplotlib} and does not install it.  The figure is
drawn by whatever CPython you already have on the machine, running as a
separate process that FEMM starts on demand and talks to over a local socket.
Nothing is sent anywhere off the machine, and nothing is installed by using
this feature.

\section{What you need}

A CPython 3.9 or newer with {\tt matplotlib}, and a window toolkit for it to
draw with.  On Windows and macOS that usually means just

\begin{verbatim}
    pip install matplotlib
\end{verbatim}

\noindent because the installers from {\tt python.org} already include Tk,
which is the toolkit {\tt matplotlib} falls back on.  Nothing else is needed.

If Tk is absent --- Homebrew's Python omits it, and most Linux distributions
package it separately --- install it alongside: {\tt brew install python-tk} on
macOS, {\tt python3-tk} on Debian and Ubuntu, {\tt py311-tkinter} on FreeBSD.
On FreeBSD use {\tt pkg} rather than {\tt pip} for {\tt matplotlib} as well,
since there are no prebuilt wheels there and {\tt pip} would try to compile it
from source.

{\tt PyQt6} or {\tt PySide6} are {\em optional}.  If one of them is installed
FEMM prefers it, giving the Qt window style instead of the Tk one; the plots
are identical either way.

FEMM finds a suitable interpreter by itself the first time a script plots
something.  If none of the candidates will do, the error says which were tried
and what was wrong with each:

\begin{verbatim}
    BridgeError: no usable Python found for the matplotlib bridge.
    Tried: python.exe (has no interactive matplotlib backend)
           python3.exe (could not be launched)
    Install it with: pip install matplotlib. The python.org installer
    already provides the Tk toolkit it needs; PyQt6 is optional and only
    changes the window style
\end{verbatim}

To point FEMM at a particular interpreter --- if you keep several, or use a
virtual environment --- set it under {\bf Edit\,$\rightarrow$\,Scripting
Preferences}, or from a script:

\begin{verbatim}
    import pybridge
    pybridge.set_interpreter("C:/Users/me/venv/Scripts/python.exe")
\end{verbatim}

The setting is remembered.  The {\tt FEMMQT\_PYTHON} environment variable does
the same thing for a single run.

\section{Which plot command to use}

FEMM has a built-in {\tt plot()} for quick $x$--$y$ curves --- it takes an
$x$ array and up to seven $y$ arrays, with optional labels, and opens a plain
plot window.  It needs nothing installed, appears instantly, and is the right
tool for a look at a sweep while you are working.

Reach for {\tt matplotlib} when you want what {\tt plot()} deliberately does
not do: contour and density plots, subplots, log axes, annotations, colour
bars, control over fonts and sizes, and vector output for a paper.

\section{Plotting a solution on its own mesh}

The example {\tt examples/plot-mesh-matplotlib.py} solves a model, reads its
finite-element mesh back, and draws the vector potential as a filled contour on
the actual triangulation rather than on a resampled grid.  The core of it is:

\begin{verbatim}
    import numpy as np, pybridge

    nn, ne = int(mo_numnodes()), int(mo_numelements())

    x = np.zeros(nn)
    y = np.zeros(nn)
    a = np.zeros(nn)
    for k in range(1, nn + 1):
        nx, ny = mo_getnode(k)
        x[k-1], y[k-1] = nx, ny
        a[k-1] = mo_getpointvalues(nx, ny)[0]

    tri = np.zeros((ne, 3), dtype=np.int32)
    for k in range(1, ne + 1):
        e = mo_getelement(k)
        tri[k-1][0] = int(e[0]) - 1     # FEMM numbers nodes from 1
        tri[k-1][1] = int(e[1]) - 1
        tri[k-1][2] = int(e[2]) - 1

    plt = pybridge.import_module("matplotlib.pyplot")
    fig, ax = plt.subplots()
    ax.tricontourf(x, y, tri, a, levels=40)
    ax.set_aspect("equal")
    plt.show()
\end{verbatim}

On a 73,000-node, 146,000-element model that reads back in about 1.4 seconds
and draws in about 0.6.

Two details in that fragment matter more than they look.

\subsection*{Keep bulk numbers in arrays, not lists}

{\tt x}, {\tt y} and {\tt a} are {\tt numpy} arrays that are allocated once and
filled in place.  The obvious alternative --- appending to a Python list --- is
much slower, and not only for the list itself.  A list of 73,000 floats is
73,000 separate objects on the interpreter's heap, and the garbage collector
walks all of them every time anything else allocates.  Measured on the mesh
above, a loop that reads 20,000 elements takes

\begin{center}
\begin{tabular}{lr}
with the node data in a {\tt numpy} array & 1.9 ms \\
with the same data in a Python list       & 69.7 ms \\
\end{tabular}
\end{center}

\noindent and the gap widens as the mesh grows.  The array also crosses to
{\tt matplotlib} as raw bytes rather than element by element, so the transfer
costs about 60 ms for eight megabytes.

\subsection*{Use {\tt int32} for indices}

The triangle array is declared {\tt dtype=np.int32}.  FEMM's {\tt numpy} is
{\tt ulab}, whose next-widest integer is {\tt int16}, and a mesh passes its
32,767 limit almost immediately --- silently, with the index wrapping to a
negative number.  {\tt int32} holds a little over two billion, which is enough
for any mesh, and it halves the bytes compared with storing indices as floats.

\section{Other packages}

The bridge is not specific to {\tt matplotlib}, and it does not insist on it
either: if the Python it finds has no plotting backend, FEMM will still use it
for everything else, and only a request for a figure fails.  Anything
importable by the CPython it found can be used:

\begin{verbatim}
    import pybridge
    csv = pybridge.import_module("csv")

    with pybridge.open("results.csv", "w", newline="") as f:
        w = csv.writer(f)
        w.writerow(["angle", "torque"])
        w.writerows(data)
\end{verbatim}

Note {\tt pybridge.open()}: the file is opened by the {\em other} process, so
that a CPython reader or writer gets a file object it can actually use.  Both
processes see the same filesystem, so the path means the same thing to both.

Spreadsheets are the other common want, and {\tt xlsxwriter} is close enough
to universal that it gets the same treatment as {\tt matplotlib}: just import
it.  Formats, charts and all.

\begin{verbatim}
    import xlsxwriter
    wb = xlsxwriter.Workbook("sweep.xlsx")
    ws = wb.add_worksheet("torque")
    bold = wb.add_format({"bold": True})

    ws.write_row(0, 0, ["angle", "torque"], bold)
    ws.write_column(1, 0, angle)          # a whole column in one call
    ws.write_column(1, 1, torque)
    wb.close()
\end{verbatim}

Note {\tt write\_column} rather than a loop over {\tt write}.  Every call is a
round trip, so writing 720 cells one at a time costs about 85 ms where two
column writes cost one --- the same rule as everywhere else here, which is to
cross in bulk rather than per item.

Archives are the third: {\tt zipfile} is forwarded too, so bundling up a run's
outputs is the code you would write anywhere.

\begin{verbatim}
    import zipfile
    with zipfile.ZipFile("run.zip", "w", zipfile.ZIP_DEFLATED) as z:
        z.write("coil.fem")
        z.write("coil.ans")
        z.writestr("notes.txt", "swept 0 to 45 degrees")
\end{verbatim}

\noindent
Note that these are {\em paths}: both processes see the same filesystem, so
the worker opens the files itself and nothing large crosses between them.

\subsection*{What crosses, and what does not}

Numbers, strings, arrays and containers of those cross by value; anything else
crosses as a proxy, and calling a method on a proxy just works.  There are two
things that cannot cross at all, and both are worth knowing before you meet
them:

\begin{itemize}
\item{\bf An open file} belongs to the process that opened it.
{\tt Workbook("sweep.xlsx")} above is fine, because that is a {\em path} ---
the worker opens the file itself.  But handing a file {\em object} to a
CPython library cannot work:

\begin{verbatim}
    f = open("sweep.xlsx", "wb")        # FEMM's file
    wb = xlsxwriter.Workbook(f)         # error: a FileIO cannot cross
\end{verbatim}

\noindent
Use \verb+pybridge.open(path, mode)+ instead, which opens the file in the
worker and hands back something CPython can write to.  The reverse direction
works as well: FEMM's own {\tt csv} module will happily write to a
{\tt pybridge.open()} file.

\item{\bf A function} defined in your script cannot be passed as a callback.
Define it in the worker with {\tt pybridge.exec()} and refer to it by name.
\end{itemize}

\noindent
Both raise immediately and say which of the two remedies applies, so this is
something you find out at the call rather than in the output.

Useful entry points, all in the {\tt pybridge} module:

\begin{itemize}
\item{\tt pybridge.import\_module("name")} --- any CPython module, as a proxy.

\item{\tt pybridge.eval("expr")}, {\tt pybridge.exec("statements")} ---
evaluate or run code in the other process.  The namespace persists between
calls and starts with {\tt np} (CPython's real {\tt numpy}) and {\tt plt}
already imported.

\item{\tt pybridge.put("name", value)} --- send an array or list across and
bind it to a name there.  This is the efficient way to move a lot of data.

\item{\tt pybridge.fetch(obj)} --- bring a value back by value; an array
returns as a {\tt numpy} array here.

\item{\tt pybridge.info()} --- which interpreter, which versions, which
plotting backend is in use.

\item{\tt pybridge.restart()}, {\tt pybridge.shutdown()} --- discard the
helper process and its windows, or stop it.  It restarts by itself on next
use.
\end{itemize}

\section{Which {\tt numpy}, which {\tt scipy}}

Inside FEMM, {\tt numpy} and {\tt scipy} are built in.  They run in the same
process, with no data copied anywhere, and they are what {\tt import numpy}
gives you --- not the CPython packages of the same name.  For almost all
post-processing they are the right answer, and reaching across the bridge for
them would be a straight loss.

They cover more than the word ``built-in'' suggests: some 115 {\tt numpy}
names, including a complete {\tt numpy.linalg} ({\tt solve}, {\tt inv},
{\tt det}, {\tt eig}, {\tt qr}, {\tt svd}, {\tt pinv}, {\tt cond},
{\tt rank}, {\tt cholesky}, {\tt norm}, all of them working on complex
matrices), {\tt numpy.fft}, {\tt polyfit} and {\tt polyval},
{\tt interp}, {\tt convolve}, {\tt trapz}, {\tt loadtxt} and {\tt savetxt};
and from {\tt scipy}, {\tt integrate.quad}, {\tt optimize.bisect},
{\tt optimize.newton}, {\tt optimize.fmin}, {\tt signal.sosfilt} and the
{\tt special} functions.

What they do not have, and where the bridge earns its keep:

\begin{itemize}
\item{\bf Curve fitting} beyond polynomials --- there is no
{\tt curve\_fit} or {\tt least\_squares}.
\item{\bf Spline interpolation} --- {\tt interp} is linear only.
\item{\bf Differential equations} --- no {\tt odeint} or {\tt solve\_ivp}.
\item{\bf Spectral estimation and filter design} --- {\tt sosfilt} applies a
filter but nothing designs one, and there is no {\tt welch} or
{\tt periodogram}.
\item{\bf Statistics} --- no {\tt scipy.stats} at all.
\end{itemize}

For those, ask for the CPython package by name:

\begin{verbatim}
    import pybridge

    pybridge.put("f", frequency)          # arrays computed in FEMM
    pybridge.put("z", impedance)
    pybridge.exec("from scipy.interpolate import CubicSpline")
    pybridge.exec("s = CubicSpline(f, z)")
    fine = pybridge.fetch(pybridge.eval("s(np.linspace(0, 1000, 2000))"))
\end{verbatim}

Note that there is no {\tt import scipy.interpolate} shim, deliberately.  The
longer spelling is a reminder that the call is leaving FEMM: each one is a
round trip to the other process, which is nothing when you call it once on a
finished array and a great deal if you call it inside a loop over elements.
Doing the arithmetic with the built-in {\tt numpy} and crossing once, at the
end, is almost always the right shape.

The other reason to keep the two straight is that a FEMM-side function cannot
cross.  This does not work:

\begin{verbatim}
    def residual(x):                      # a function defined here
        return measured - model(x)
    opt.least_squares(residual, x0)       # BridgeError
\end{verbatim}

\noindent because the optimiser would have to call back into FEMM.  Define the
function on the far side with {\tt pybridge.exec} and it works normally; or
keep the loop here, using the built-in {\tt scipy.optimize.fmin} or
{\tt newton}, and call FEMM from inside it.  The second is usually what you
want anyway when each evaluation is a finite-element solve.

\section{Limitations}

\begin{itemize}
\item{\bf Functions cannot cross.}  A FEMM-side function cannot be passed to
CPython, so there is no {\tt key=} argument to {\tt sorted}, no callback to
{\tt FuncAnimation}, and no interactive widget or pick event.  Define such a
function in the other process with {\tt pybridge.exec} instead.

\item{\bf {\tt numpy} and {\tt scipy} stay local.}  Those names refer to
FEMM's own built-in array library, which runs in-process and is much faster
for ordinary work than sending each operation across.  When you specifically
want CPython's {\tt numpy}, ask for it: {\tt pybridge.import\_module("numpy")},
or use the {\tt np} already present in the {\tt eval} namespace.

\item{\bf Figures belong to the other process.}  They appear in their own
windows and cannot be embedded in a FEMM window.

\item{\bf One call at a time.}  While a plotting call is in progress FEMM will
refuse a second one rather than let two overlap.  This can only arise if a
script is somehow re-entered while it waits.

\item Pressing {\bf Esc} during a plotting call that has hung will abandon it.
The helper process is discarded, and the next plotting command starts a fresh
one; any figures it was holding are lost.
\end{itemize}
